\documentclass[runningheads]{llncs}

% ---------------------------------------------------------------
% ECCV package must load first (provides xcolor, url, etc.)
\usepackage[review,year=2026,ID=*****]{eccv}

% ---------------------------------------------------------------
% Other packages
\usepackage{eccvabbrv}
\usepackage{graphicx}
\usepackage{booktabs}
\usepackage{amsmath,amssymb}
\usepackage{multirow}
\usepackage{listings}
\usepackage{algorithm}
\usepackage{algorithmic}
\usepackage[accsupp]{axessibility}
\usepackage[pagebackref,breaklinks,colorlinks,citecolor=eccvblue]{hyperref}
\usepackage[capitalize]{cleveref}
\usepackage{orcidlink}

% Listing style for prompt templates
\lstset{
  basicstyle=\ttfamily\footnotesize,
  breaklines=true,
  frame=single,
  backgroundcolor=\color{gray!5},
  rulecolor=\color{gray!40},
  xleftmargin=4pt,
  xrightmargin=4pt,
  aboveskip=6pt,
  belowskip=6pt,
}

\begin{document}

\title{UrbanAlign: Supplementary Material}
\titlerunning{UrbanAlign --- Supplementary}
\author{Anonymous ECCV 2026 Submission}
\authorrunning{Anonymous}
\institute{Anonymous Institution}
\maketitle

This supplementary document provides additional details that support the main paper.
\Cref{sec:dataset} describes the dataset filtering protocol and data-split statistics.
\Cref{sec:lwrr_weights} presents per-pair LWRR weight variation analysis.
\Cref{sec:sensitivity_full} reports full parameter-sweep grids expanding on the main-text sensitivity analysis.
\Cref{sec:prompts} reproduces the exact prompt templates used for multi-agent feature distillation.
\Cref{sec:e2e} details the end-to-end dimension optimization procedure.
\Cref{sec:vrm_gain} reports per-mode LWRR alignment gains.
\Cref{sec:algorithm} provides the complete algorithm pseudocode.


% ==================================================================
% A. Dataset & Filtering Protocol
% ==================================================================
\section{Dataset and Filtering Protocol}
\label{sec:dataset}

\subsection{Place Pulse 2.0 Overview}

Place Pulse~2.0~\cite{salesses2013collaborative} contains 1,169,078 pairwise comparisons over 110,688 Google Street View images from 56 cities worldwide.
Each comparison asks: ``Which place looks more \texttt{<category>}?'' for six perception categories: \emph{safety}, \emph{lively}, \emph{beautiful}, \emph{wealthy}, \emph{depressing}, and \emph{boring}.
Each comparison yields one of three outcomes: \emph{left}, \emph{right}, or \emph{equal}.

\subsection{Consensus Filtering ($N \geq 3$)}

Raw Place Pulse data contains many image pairs with only one or two votes, which carry high individual noise.
We apply a consensus threshold: only pairs with $N \geq 3$ independent votes and a strict majority ($> 50\%$ agreement on the winning label) are retained.
\Cref{tab:consensus} reports the filtering statistics for the \emph{wealthy} category.

\begin{table}[ht]
\centering
\caption{Consensus filtering statistics for the \emph{wealthy} category.}
\label{tab:consensus}
\small
\begin{tabular}{lcc}
\toprule
Criterion & Count & \% of Raw \\
\midrule
Total raw pairs (6 categories) & 1643 & 100\% \\
Pairs with $N \geq 3$ votes & 1643 & 100\% \\
\quad of which strict majority & 1588 & 96.7\% \\
\quad label = left & 788 & --- \\
\quad label = right & 721 & --- \\
\quad label = equal & 79 & 4.8\% \\
\midrule
Retained after filtering & 1643 & 100\% \\
\bottomrule
\end{tabular}
\end{table}

\subsection{TrueSkill Rating Computation}

We convert the filtered pairwise votes into continuous per-image ratings using TrueSkill~\cite{herbrich2006trueskill} with priors
($\mu_0{=}25$, $\sigma_0{=}8.33$, $\beta{=}4.17$, draw prob.\ $=0.10$).
After convergence, images are ranked by $\mu$; the resulting values serve two purposes:
(1)~\emph{Stage~1 consensus sampling}: selecting high-$\mu$ and low-$\mu$ exemplars for dimension extraction;
(2)~\emph{Stage~3 supervision}: providing continuous TrueSkill score differences $y^{\mathrm{TS}}_{ij} {=} \mu_i {-} \mu_j$ as regression targets for LWRR.

\subsection{Data Split Protocol}

From the filtered pairs, we randomly sample 288 pairs (sampling ratio 0.01 of the full filtered pool).
These 288 pairs are then split into:
\begin{itemize}
    \item \textbf{Reference set} $\mathcal{D}_{\mathrm{ref}}$ (70\%, ${\approx}201$ pairs): provides human labels and TrueSkill supervision for LWRR calibration.
    Mirror augmentation doubles the effective reference size to ${\approx}402$ points.
    \item \textbf{Pool set} $\mathcal{D}_{\mathrm{pool}}$ (30\%, ${\approx}87$ pairs): calibrated by Stage~3 and evaluated against held-out human labels.
\end{itemize}
The split uses a fixed random seed (42) for reproducibility.
$\mathcal{D}_{\mathrm{ref}} \cap \mathcal{D}_{\mathrm{pool}} = \varnothing$ by construction.

After Stage~3 alignment with confidence-based filtering (selection ratio 0.15 applied to pool), 87 pairs remain.
Excluding human ``equal'' labels yields 83 evaluation pairs.


% ==================================================================
% B. LWRR Per-Pair Weight Analysis
% ==================================================================
\section{LWRR Per-Pair Weight Analysis}
\label{sec:lwrr_weights}

A key feature of LWRR is that each pool pair receives its \emph{own} locally-fitted weight vector $\hat{\mathbf{w}} \in \mathbb{R}^{7}$, reflecting the dimension importances in its particular neighbourhood of the hybrid manifold.
This contrasts with a global linear model where all pairs share the same weights.

\subsection{Weight Variation Across the Manifold}

\Cref{tab:weight_variation} reports the mean and standard deviation of each dimension's local weight across all 87 aligned pool pairs.
High standard deviation indicates that the dimension's importance varies substantially across different regions of the perception manifold---precisely the heterogeneity that motivates local rather than global calibration.

\begin{table}[ht]
\centering
\caption{LWRR local weight statistics across 87 aligned pool pairs (\emph{wealthy} category, Mode~4).}
\label{tab:weight_variation}
\small
\begin{tabular}{lccc}
\toprule
Dimension & Mean $w$ & Std $w$ & CV \\
\midrule
Vegetation Maintenance    & 3.384 & 4.636 & 1.37 \\
Street Cleanliness        & $-$1.566 & 3.306 & 2.11 \\
Vehicle Quality           & $-$1.184 & 5.368 & 4.53 \\
Building Modernity        & 1.165 & 3.871 & 3.32 \\
Infrastructure Condition  & 1.006 & 3.369 & 3.35 \\
Fa\c{c}ade Quality       & $-$0.940 & 2.507 & 2.67 \\
Pavement Integrity        & $-$0.464 & 4.363 & 9.40 \\
\bottomrule
\end{tabular}
\\[4pt]
\footnotesize{CV = Coefficient of Variation ($\mathrm{Std}/|\mathrm{Mean}|$). Higher CV indicates greater spatial heterogeneity of dimension importance. Sorted by $|\text{Mean}|$.}
\end{table}

\Cref{tab:weight_variation_all} extends this analysis to all six categories.

\begin{table}[ht]
\centering
\caption{Top-3 LWRR local weights by $|\text{Mean}|$ for each perception category (Mode~4). High Std relative to Mean confirms that dimension importance varies substantially across the manifold.}
\label{tab:weight_variation_all}
\small
\setlength{\tabcolsep}{4pt}
\begin{tabular}{llccc}
\toprule
Category & Dimension & Mean $w$ & Std $w$ & CV \\
\midrule
\multirow{3}{*}{Safety ($n{=}87$)}
& Visibility Clarity      & 2.638  & 5.823 & 2.21 \\
& Signage \& Wayfinding   & 2.420  & 4.192 & 1.73 \\
& Noise Levels            & $-$1.337 & 3.255 & 2.43 \\
\midrule
\multirow{3}{*}{Beautiful ($n{=}74$)}
& Street Cleanliness      & $-$2.873 & 4.956 & 1.73 \\
& Color Coordination      & 2.473  & 4.283 & 1.73 \\
& Natural Elements        & 2.220  & 4.837 & 2.18 \\
\midrule
\multirow{3}{*}{Lively ($n{=}87$)}
& Commercial Density      & $-$0.932 & 2.193 & 2.35 \\
& Street Furniture        & $-$0.735 & 1.715 & 2.33 \\
& Building Fa\c{c}ade     & 0.688  & 2.266 & 3.29 \\
\midrule
\multirow{3}{*}{Wealthy ($n{=}87$)}
& Vegetation Maint.       & 3.384  & 4.636 & 1.37 \\
& Street Cleanliness      & $-$1.566 & 3.306 & 2.11 \\
& Vehicle Quality         & $-$1.184 & 5.368 & 4.53 \\
\midrule
\multirow{3}{*}{Boring ($n{=}80$)}
& Functional Diversity    & 1.353  & 2.661 & 1.97 \\
& Spatial Enclosure       & $-$1.000 & 3.321 & 3.32 \\
& Lighting Quality        & $-$0.838 & 2.489 & 2.97 \\
\midrule
\multirow{3}{*}{Depressing ($n{=}78$)}
& Vegetation \& Greenery  & $-$2.354 & 4.479 & 1.90 \\
& Lighting Quality        & 2.167  & 4.366 & 2.01 \\
& Traffic \& Noise        & $-$1.612 & 3.773 & 2.34 \\
\bottomrule
\end{tabular}
\end{table}

\subsection{Illustrative Examples}

We select three representative pool pairs from different regions of the manifold to illustrate how LWRR adapts its weight profile.

\medskip\noindent\textbf{Example 1: Vegetation-driven pair.}
In the suburban portion of the manifold, LWRR assigns dominant weights to Vegetation Maintenance ($w{=}3.38$) and Building Modernity ($w{=}1.17$), while Fa\c{c}ade Quality receives near-zero weight ($w{\approx}-0.94$).
This reflects the perceptual pattern that in residential areas, landscaping quality is the primary cue for perceived wealth.

\medskip\noindent\textbf{Example 2: Infrastructure-driven pair.}
In the dense urban core, the local weight profile shifts: Infrastructure Condition ($w{=}1.01$) and Vehicle Quality ($w{=}-1.18$) dominate, while Vegetation drops.
The sign reversal on Vehicle Quality indicates contextual inversion---in downtown areas, the presence of high-end vehicles may signal commercial rather than residential wealth.

\medskip\noindent\textbf{Example 3: Balanced transitional pair.}
Pairs in mixed-use zones show relatively uniform weight magnitude across all seven dimensions (std of $|w|$ across dims $<$0.5), confirming that these regions lack a single dominant perceptual cue.
The high CV (Coefficient of Variation $>$3) for Pavement Integrity (9.40) in \cref{tab:weight_variation} is driven precisely by these transitional pairs.

\medskip
These examples confirm the theoretical motivation for local calibration: the dimensions that matter most for perceiving ``wealth'' depend on the visual context.
A single global weight vector would average over these distinct regimes, losing the nuance that LWRR captures.


% ==================================================================
% C. Detailed Sensitivity Analysis
% ==================================================================
\section{Detailed Sensitivity Analysis}
\label{sec:sensitivity_full}

The main text (\cref{tab:sensitivity} in the main paper) summarises the per-category sensitivity of Stage~3 to key hyperparameters.
Here we report the full parameter-sweep grids, varying one parameter at a time while holding others at their defaults ($K{=}20$, $\alpha{=}0.3$, $\tau_{\mathrm{kernel}}{=}1.0$, $\lambda{=}1.0$, $\varepsilon{=}0.8$, $\theta{=}0.6$).
All results are Mode~4, accuracy excluding ``equal'' human labels.

\subsection{Neighbourhood Size $K_{\max}$}

\begin{table}[ht]
\centering
\caption{Full $K_{\max}$ sweep (Mode~4, all categories, accuracy \% excl.\ equal).}
\label{tab:k_sweep}
\small
\setlength{\tabcolsep}{4pt}
\begin{tabular}{ccccccc|c}
\toprule
$K_{\max}$ & Safety & Beaut. & Lively & Wealthy & Boring & Depress. & Avg \\
\midrule
10 & 64.6 & \textbf{59.1} & 45.0 & 57.5 & 52.9 & 42.0 & 53.5 \\
20 & 66.3 & 56.5 & \textbf{57.1} & \textbf{72.3} & \textbf{61.5} & 54.3 & \textbf{61.3} \\
30 & 63.7 & 52.2 & 56.9 & 66.7 & 59.1 & \textbf{55.1} & 58.9 \\
50 & \textbf{67.5} & 56.2 & 53.4 & 67.9 & 51.6 & 54.0 & 58.4 \\
\bottomrule
\end{tabular}
\end{table}

\noindent\textbf{Analysis.}
Small $K$ ($=10$) leads to high-variance local fits (each ridge regression uses very few reference points).
$K{=}20$ achieves the best average (61.3\%) and is optimal for four of six categories.
However, the optimal $K$ varies: \emph{safety} favours $K{=}50$ (67.5\%), \emph{beautiful} favours $K{=}10$ (59.1\%), reflecting different local densities of the reference manifold across categories.

\subsection{Hybrid Weight $\alpha$}

\begin{table}[ht]
\centering
\caption{Full $\alpha$ sweep (Mode~4, all categories, accuracy \% excl.\ equal).}
\label{tab:alpha_sweep}
\small
\setlength{\tabcolsep}{4pt}
\begin{tabular}{cccccccc|c}
\toprule
$\alpha$ & Sem.\ Wt. & Safety & Beaut. & Lively & Wealthy & Boring & Depress. & Avg \\
\midrule
0.1 & 90\% & 65.4 & 57.6 & 51.3 & 67.5 & 57.4 & 55.1 & 59.1 \\
0.2 & 80\% & 63.7 & 50.0 & 56.4 & 66.7 & 52.3 & 57.7 & 57.8 \\
0.3 & 70\% & \textbf{66.3} & 56.5 & \textbf{57.1} & \textbf{72.3} & \textbf{61.5} & 54.3 & \textbf{61.3} \\
0.5 & 50\% & 65.4 & \textbf{63.3} & 52.7 & 65.8 & 54.3 & \textbf{60.0} & 60.3 \\
0.7 & 30\% & 64.1 & 59.1 & 56.8 & 64.2 & 50.7 & 45.6 & 56.7 \\
\bottomrule
\end{tabular}
\end{table}

\noindent\textbf{Analysis.}
$\alpha{=}0.3$ (70\% semantic weight) is optimal on average, but per-category optima differ notably: \emph{beautiful} favours $\alpha{=}0.5$ (63.3\%) and \emph{depressing} also favours $\alpha{=}0.5$ (60.0\%), suggesting that CLIP features carry more discriminative signal for aesthetics-related categories.
Pure CLIP ($\alpha{=}1.0$) performs poorly because CLIP's pre-training distribution (web images of objects) is mismatched with street-view perception tasks.
This empirically validates our hybrid design while highlighting category-specific tuning opportunities.

\subsection{Kernel Bandwidth $\tau_{\mathrm{kernel}}$}

\begin{table}[ht]
\centering
\caption{Full $\tau_{\mathrm{kernel}}$ sweep (Mode~4, \emph{wealthy}).}
\label{tab:tau_sweep}
\small
\begin{tabular}{cccc}
\toprule
$\tau_{\mathrm{kernel}}$ & Acc (\%) & $\kappa$ & $N_{\mathrm{aligned}}$ \\
\midrule
0.1 & 67.9 & 0.35 & 87 \\
0.3 & 70.4 & 0.40 & 87 \\
0.5 & 70.7 & 0.41 & 87 \\
0.8 & 72.0 & 0.43 & 87 \\
1.0 & 72.3 & 0.44 & 87 \\
\bottomrule
\end{tabular}
\end{table}

\noindent\textbf{Analysis.}
Smaller $\tau$ concentrates weight on the nearest neighbours (sharper kernel); larger $\tau$ approaches uniform weighting within the $K$-neighbourhood.
The relatively flat response across the tested range ($<$4.4\,pp variation) suggests that the $K$-NN pre-selection already isolates relevant neighbours, making the exact kernel bandwidth less critical.

\subsection{Ridge Penalty $\lambda_{\mathrm{ridge}}$}

\begin{table}[ht]
\centering
\caption{Full $\lambda_{\mathrm{ridge}}$ sweep (Mode~4, \emph{wealthy}).}
\label{tab:lambda_sweep}
\small
\begin{tabular}{cccc}
\toprule
$\lambda$ & Acc (\%) & $\kappa$ & $N_{\mathrm{aligned}}$ \\
\midrule
0.01 & 70.7 & 0.41 & 87 \\
0.05 & 70.7 & 0.41 & 87 \\
0.1  & 70.7 & 0.41 & 87 \\
0.5  & 70.7 & 0.41 & 87 \\
1.0  & 72.3 & 0.44 & 87 \\
\bottomrule
\end{tabular}
\end{table}

\noindent\textbf{Analysis.}
The ridge penalty $\lambda$ has minimal impact across several orders of magnitude ($<$1.6\,pp variation).
This is expected: with $K{=}20$ neighbours and $n{=}7$ dimensions, the system is well-overdetermined ($K \gg n$), so regularisation has little effect.

\subsection{Equal Threshold $\varepsilon$ (EQUAL\_EPS)}

\begin{table}[ht]
\centering
\caption{Full $\varepsilon$ (EQUAL\_EPS) sweep (Mode~4, all categories, accuracy \% excl.\ equal).}
\label{tab:eps_sweep}
\small
\setlength{\tabcolsep}{4pt}
\begin{tabular}{ccccccc|c}
\toprule
$\varepsilon$ & Safety & Beaut. & Lively & Wealthy & Boring & Depress. & Avg \\
\midrule
0.3 & 66.3 & 56.5 & 57.5 & 72.3 & 59.7 & 52.8 & 60.9 \\
0.5 & 66.3 & 56.5 & 56.4 & 72.3 & 59.4 & 52.8 & 60.6 \\
0.8 & 66.3 & 56.5 & 57.1 & 72.3 & \textbf{61.5} & 54.3 & 61.3 \\
1.2 & 66.2 & 57.4 & 56.6 & \textbf{72.8} & \textbf{63.5} & 54.4 & 61.8 \\
2.0 & \textbf{67.1} & \textbf{59.7} & \textbf{59.1} & 72.4 & 63.3 & \textbf{56.2} & \textbf{63.0} \\
\bottomrule
\end{tabular}
\end{table}

\noindent\textbf{Analysis.}
The equal threshold $\varepsilon$ shows a mild but consistent trend: larger $\varepsilon$ (more permissive equal classification) improves accuracy across most categories, with the average rising from 60.9\% ($\varepsilon{=}0.3$) to 63.0\% ($\varepsilon{=}2.0$).
This suggests that many borderline pairs are better classified as non-equal, potentially because the VLM's score differences carry marginal but genuine signal even for close pairs.

\subsection{Combined Search}

In addition to the one-at-a-time sweeps above, we performed a random search over the combined parameter space ($N{=}50{,}000$ configurations) to check for interactions.
\Cref{tab:combined_sweep} reports the per-category best configurations found.

\begin{table}[ht]
\centering
\caption{Best per-category configurations from combined random search (Mode~4, accuracy \% excl.\ equal).}
\label{tab:combined_sweep}
\small
\setlength{\tabcolsep}{4pt}
\begin{tabular}{lccl}
\toprule
Category & Acc.\ (\%) & $\kappa$ & Key Parameters \\
\midrule
Safety      & 86.0 & 0.703 & $K{=}30$, $\alpha{=}0.1$, sel${=}0.6$ \\
Beautiful   & 70.7 & 0.407 & $K{=}30$, $\alpha{=}0.5$, sel${=}0.6$ \\
Lively      & 72.3 & 0.446 & $K{=}30$, $\alpha{=}0.3$, sel${=}0.6$ \\
Wealthy     & 80.0 & 0.599 & $K{=}30$, $\alpha{=}0.1$, sel${=}0.6$ \\
Boring      & 69.8 & 0.402 & $K{=}30$, $\alpha{=}0.3$, sel${=}0.6$ \\
Depressing  & 66.7 & 0.333 & $K{=}30$, $\alpha{=}0.5$, sel${=}0.6$ \\
\midrule
Average     & 74.2 & 0.482 & --- \\
\bottomrule
\end{tabular}
\end{table}

\noindent The combined search yields 74.2\% average accuracy ($\kappa{=}0.482$), a +12.9\,pp gain over the default configuration (61.3\%), demonstrating significant headroom from per-category hyperparameter tuning.
The most consistent finding is $K{=}30$ and $\text{sel}{=}0.6$ across all categories, while the optimal $\alpha$ varies (0.1 for safety/wealthy, 0.3 for lively/boring, 0.5 for beautiful/depressing).


% ==================================================================
% D. Prompt Templates
% ==================================================================
\section{Prompt Templates}
\label{sec:prompts}

This section reproduces the exact prompt templates used in Stage~2 multi-agent feature distillation (Mode~4: pairwise + multi-agent).
All three agents receive the image pair and the Stage~1--extracted dimension definitions as input.
Temperature values: Observer $\tau{=}0.3$, Debater $\tau{=}0.5$, Judge $\tau{=}0.1$.

\subsection{Stage 1: Dimension Extraction Prompt}

\begin{lstlisting}[caption={Dimension extraction prompt (Stage 1).}]
You are an expert urban perception researcher. I will show
you street-view images that have been rated by crowdsourced
participants on the perception dimension: "{category}".

HIGH-rated images (TrueSkill mu > 75th percentile):
{high_image_descriptions}

LOW-rated images (TrueSkill mu < 25th percentile):
{low_image_descriptions}

Based on the visual differences between HIGH and LOW images,
identify 5-10 specific, visually observable dimensions that
explain why some scenes are perceived as more "{category}"
than others.

Output JSON format:
{
  "dimensions": [
    {
      "name": "Dimension Name",
      "description": "What this dimension measures",
      "high_indicator": "Visual cues for high scores",
      "low_indicator": "Visual cues for low scores"
    }
  ]
}

Requirements:
- Each dimension must be visually observable from a street
  view image
- Each dimension must be continuously scorable (1-10 scale)
- Dimensions should be universally applicable across cities
\end{lstlisting}

\subsection{Observer Prompt}

\begin{lstlisting}[caption={Observer agent prompt (Stage 2, Mode 4).}]
You are an objective Observer. Examine two street-view
images and describe what you see along each evaluation
dimension. Do NOT make judgments or comparisons -- only
describe observable visual details.

Category: {category}
Dimensions:
{dimension_definitions}

For each dimension, describe what you observe in:
- Image A: [visual details]
- Image B: [visual details]

Output JSON:
{
  "observations": {
    "dimension_name": {
      "image_a": "description of what you see",
      "image_b": "description of what you see"
    }
  }
}
\end{lstlisting}

\subsection{Debater Prompt}

\begin{lstlisting}[caption={Debater agent prompt (Stage 2, Mode 4).}]
You are a Debater. Given the Observer's descriptions of
two street-view images, argue BOTH sides for each
dimension -- why Image A might score higher AND why
Image B might score higher.

Observer's descriptions:
{observer_output}

Category: {category}
Dimensions:
{dimension_definitions}

For each dimension, provide:
- Argument for Image A scoring higher
- Argument for Image B scoring higher
- Key uncertainties or ambiguities

Output JSON:
{
  "debates": {
    "dimension_name": {
      "argument_for_a": "why A might score higher",
      "argument_for_b": "why B might score higher",
      "uncertainties": "what is ambiguous"
    }
  }
}
\end{lstlisting}

\subsection{Judge Prompt}

\begin{lstlisting}[caption={Judge agent prompt (Stage 2, Mode 4).}]
You are the final Judge. Given the Observer's descriptions
and Debater's arguments, produce final scores for both
images on each dimension.

Observer's descriptions:
{observer_output}

Debater's arguments:
{debater_output}

Category: {category}
Dimensions:
{dimension_definitions}

Score each image on each dimension from 1 (lowest) to 10
(highest). Also determine the overall winner.

Output JSON:
{
  "image_a_scores": {
    "dimension_name": score (1-10)
  },
  "image_b_scores": {
    "dimension_name": score (1-10)
  },
  "overall_intensity_a": weighted_sum,
  "overall_intensity_b": weighted_sum,
  "winner": "left" | "right" | "equal",
  "ai_dimension_weights": {
    "dimension_name": weight (0-1)
  }
}
\end{lstlisting}

% ==================================================================
% E. End-to-End Dimension Optimization
% ==================================================================
\section{End-to-End Dimension Optimization}
\label{sec:e2e}

The semantic dimensions discovered in Stage~1 depend on the VLM's sampling temperature and the particular consensus exemplars presented.
A natural question is: \emph{can we automatically find the dimension set that maximises end-to-end alignment accuracy?}

\smallskip\noindent\textbf{Formulation.}
Let $\mathcal{D}_c^{(t)}$ denote the dimension set generated at trial $t$ with generation temperature $\tau_{\mathrm{gen}}^{(t)}$.
Each trial runs the full pipeline: dimension extraction (Stage~1) $\to$ pairwise scoring (Stage~2, Mode~2 for cost efficiency) $\to$ LWRR alignment (Stage~3) $\to$ accuracy evaluation.
The optimal dimension set is:
\begin{equation}
\mathcal{D}_c^{*} = \arg\max_{t \in \{1,\ldots,T\}} \;\mathrm{Acc}\bigl(\text{LWRR}\bigl(\text{Score}(\mathcal{D}_{\mathrm{pool}},\, \mathcal{D}_c^{(t)}),\, \mathcal{D}_{\mathrm{ref}}\bigr)\bigr).
\label{eq:e2e_supp}
\end{equation}

\smallskip\noindent\textbf{Search strategy.}
Rather than exhaustive search over the combinatorial space of possible dimensions, we exploit the VLM's generative diversity by progressively increasing the sampling temperature across trials: $\tau_{\mathrm{gen}}^{(t)} = 0.7 + 0.1t$ for $t=0,\ldots,T{-}1$.
Low temperatures produce conservative, high-consensus dimensions; higher temperatures encourage more creative, potentially complementary dimensions.
Each trial uses a fixed small evaluation sample (${\sim}150$ pairs) shared across all trials to ensure comparability.

\smallskip\noindent\textbf{Cost analysis.}
Each trial requires one Stage~1 call (${\sim}1$ API call) and ${\sim}150$ Stage~2 scoring calls.
With $T$ trials, total cost is ${\sim}T \times 151$ API calls.
For $T{=}8$, this amounts to ${\approx}1{,}200$ calls---a modest overhead given that Stage~3 LWRR alignment is purely local computation with zero API cost.

\begin{table}[ht]
\centering
\caption{End-to-end dimension optimization results (\emph{wealthy} category, $T{=}6$ trials). Each trial generates a new dimension set at increasing temperature, scores a fixed evaluation sample via Mode~2, and measures post-LWRR alignment accuracy.}
\label{tab:e2e}
\small
\begin{tabular}{ccccp{5.8cm}}
\toprule
Trial & $\tau_{\mathrm{gen}}$ & Acc.\,(\%) & $\kappa$ & Key Dimensions \\
\midrule
\textbf{0} & \textbf{0.7} & \textbf{56.3} & \textbf{0.231} & Fa\c{c}ade Quality, Vegetation Maint., Pavement Integrity, Vehicle Quality, Building Modernity, Infrastructure, Cleanliness, Street Lighting, Public Art, + 8 more (17 total) \\
1 & 0.8 & 50.6 & 0.100 & 10 dimensions \\
2 & 0.9 & 52.9 & 0.148 & 10 dimensions \\
3 & 1.0 & 49.4 & 0.089 & 10 dimensions \\
4 & 1.1 & 48.3 & 0.081 & 10 dimensions \\
5 & 1.2 & 51.7 & 0.127 & 10 dimensions \\
\midrule
Best & 0.7 & 56.3 & 0.231 & Trial 0 (17 dims) \\
\bottomrule
\end{tabular}
\\[4pt]
\footnotesize{$\tau_{\mathrm{gen}}$: VLM generation temperature for dimension extraction. Evaluation sample: ${\sim}87$ pool pairs.
Note: E2E trials use Mode~2 (single-shot pairwise) for cost efficiency; main results use Mode~4.}
\end{table}

\smallskip\noindent\textbf{Results.}
The lowest-temperature trial ($\tau_{\mathrm{gen}}{=}0.7$, Trial~0) achieves the best performance (56.3\% accuracy, $\kappa{=}0.231$), with performance degrading monotonically at higher temperatures.
Trial~0 generates a richer dimension set (17~dimensions vs.\ the typical 10) including domain-specific additions such as \emph{Street Lighting}, \emph{Public Art Presence}, \emph{Signage Quality}, \emph{Street Furniture Quality}, \emph{Traffic Flow}, \emph{Noise Levels}, \emph{Pedestrian Accessibility}, \emph{Color Palette}, \emph{Historical Features}, and \emph{Cultural Elements}.
This finding contradicts the prior expectation that moderate temperatures would be optimal: instead, the VLM's most conservative, consensus-aligned output is most effective.
The monotonic degradation with temperature confirms the $(1{-}p^*)^T$ convergence bound (\cref{eq:convergence_supp}): the optimal dimension set has highest generation probability at low $\tau$, and increasing $\tau$ reduces this probability.

Note that the E2E trials use Mode~2 (single-shot pairwise) for cost efficiency rather than Mode~4 (multi-agent pairwise); the absolute accuracy values are therefore lower than the Mode~4 results reported in the main text.
The key finding is the \emph{ranking} of dimension sets, which transfers across scoring modes.

\subsection{Theoretical Analysis}

\noindent\textbf{Information-bottleneck interpretation.}
Each dimension set $\mathcal{D}_c^{(t)}$ defines a different compression of the visual input through the lens of the information bottleneck~\cite{tishby2000information}.
Let $X$ denote the raw visual input and $Y$ the human label.
Stage~2 compresses $X$ into dimension scores $S^{(t)}=f_{\mathcal{D}_c^{(t)}}(X)$, and Stage~3 maps $S^{(t)}$ to prediction $\hat{Y}$.
The optimal dimension set maximises $I(\hat{Y}; Y)$ subject to the constraint that the bottleneck representation $S^{(t)}$ remains compact ($n \in [5,10]$ dimensions).
Equation~\eqref{eq:e2e_supp} selects the compression that maximises this downstream mutual information empirically.

\smallskip\noindent\textbf{Convergence guarantee.}
Let $\mathcal{F}$ denote the (finite) set of dimension sets reachable by the VLM at any temperature in $[\tau_{\min}, \tau_{\max}]$.
Although $|\mathcal{F}|$ is large (combinatorially many possible dimension names and descriptions), the VLM's generative distribution concentrates probability mass on a much smaller effective set $\mathcal{F}_{\mathrm{eff}} \subset \mathcal{F}$ of semantically coherent dimension sets.
If we model each trial as drawing independently from the VLM's distribution over $\mathcal{F}_{\mathrm{eff}}$, then after $T$ trials the probability of missing the optimal set $\mathcal{D}_c^*$ satisfies:
\begin{equation}
\Pr\bigl[\mathcal{D}_c^* \notin \{\mathcal{D}_c^{(1)}, \ldots, \mathcal{D}_c^{(T)}\}\bigr] \leq (1 - p^*)^T,
\label{eq:convergence_supp}
\end{equation}
where $p^* = \Pr_{\mathrm{VLM}}[\mathcal{D}_c^{(t)} = \mathcal{D}_c^*]$ is the probability of generating the optimal set in a single trial.
For a near-optimal set (within $\epsilon$ accuracy of the best), this probability is typically higher.
Concretely, if the top-$k$ dimension sets are each generated with probability $\geq p_{\min}$, then $T \geq \lceil \ln(1/\delta) / \ln(1/(1-k \cdot p_{\min})) \rceil$ trials suffice to find at least one with probability $\geq 1-\delta$.

\smallskip\noindent\textbf{Search strategy: progressive temperature as exploration--exploitation.}
The progressive temperature schedule $\tau_{\mathrm{gen}}^{(t)} = 0.7 + 0.1t$ implements a structured exploration of the dimension space that balances two objectives:
\begin{itemize}
    \item \emph{Exploitation} (low $\tau$): conservative temperatures ($\tau \leq 0.8$) produce high-probability, consensus-aligned dimensions---the VLM's ``default'' understanding of the perception category.
    \item \emph{Exploration} (high $\tau$): elevated temperatures ($\tau \geq 1.0$) increase sampling entropy, generating creative or less conventional dimensions that may capture perceptual nuances missed by conservative sampling.
\end{itemize}
This is analogous to random search with a guided prior~\cite{bergstra2012random}: rather than sampling uniformly from the combinatorial space of possible dimensions, we leverage the VLM's own generative distribution as an implicit prior over useful dimension sets.
The temperature schedule ensures that this prior is traversed from its mode (exploitation) toward its tails (exploration).
Unlike Bayesian optimisation~\cite{snoek2012practical}, which would require a surrogate model over the discrete, variable-length dimension space, our approach avoids the need for a surrogate entirely.

\smallskip\noindent\textbf{Sample complexity for evaluation.}
Each trial evaluates alignment accuracy on a fixed sample of $m$ pairs.
By Hoeffding's inequality, the empirical accuracy $\widehat{\mathrm{Acc}}$ satisfies $|\widehat{\mathrm{Acc}} - \mathrm{Acc}| \leq \sqrt{\ln(2T/\delta)/(2m)}$ uniformly over all $T$ trials with probability $\geq 1-\delta$.
For $T{=}8$ and $m{=}150$, this gives a confidence interval of $\pm 8.5$\% at $\delta{=}0.05$---sufficient to distinguish meaningfully different dimension sets (where accuracy gaps are typically $>$10\,pp) while keeping API costs modest.
Larger evaluation samples would tighten this bound but at proportionally higher cost.

\smallskip\noindent\textbf{Connection to LWRR convergence.}
The end-to-end accuracy depends not only on the dimension set but also on the LWRR calibration quality.
For regularised least-squares with $K$ local neighbours in $n$ dimensions, the excess risk decays as $\mathcal{O}(\lambda + n/(K\lambda))$~\cite{caponnetto2007optimal,hoerl1970ridge}, which is minimised at $\lambda^* \propto (n/K)^{1/2}$.
In our setting ($n{=}7$, $K{=}20$), the system is well-overdetermined ($K \gg n$), so the LWRR calibration error is small regardless of the dimension set---confirming that the E2E accuracy variation across trials is driven primarily by the quality of the dimension-level representation, not by the calibration algorithm.


% ==================================================================
% F. LWRR Alignment Gain Per Mode
% ==================================================================
\section{LWRR Alignment Gain Per Mode}
\label{sec:vrm_gain}

\Cref{tab:vrm_gain} reports the accuracy change from raw VLM output to post-LWRR calibration for each of the four factorial modes on \emph{wealthy}.
\Cref{tab:vrm_gain_allcat} extends the Mode~4 analysis to all six categories.

\begin{table}[ht]
\centering
\caption{LWRR alignment gain (accuracy \% excluding ``equal'') per mode, \emph{wealthy} category. Only Mode~4 benefits from calibration.}
\label{tab:vrm_gain}
\small
\begin{tabular}{lcccc}
\toprule
 & Mode~1 & Mode~2 & Mode~3 & Mode~4 \\
\midrule
Raw Acc.\ (\%)     & 69.1 & 64.2 & 65.7 & 65.0 \\
Aligned Acc.\ (\%) & 59.0 & 64.2 & 54.1 & 72.3 \\
$\Delta$ (pp)      & $-$10.1 & $\pm$0.0 & $-$11.6 & \textbf{+7.3} \\
\midrule
Raw $\kappa$       & 0.383 & 0.279 & 0.314 & 0.298 \\
Aligned $\kappa$   & 0.179 & 0.284 & 0.082 & \textbf{0.440} \\
$n$ (pool pairs)   & 78 & 81 & 74 & 83 \\
\bottomrule
\end{tabular}
\end{table}

\begin{table}[ht]
\centering
\caption{LWRR alignment gain for Mode~4 across all six perception categories (accuracy \% excluding ``equal''). LWRR improves four of six categories.}
\label{tab:vrm_gain_allcat}
\small
\setlength{\tabcolsep}{3.5pt}
\begin{tabular}{lcccccc|c}
\toprule
 & Safety & Beaut. & Lively & Wealthy & Boring & Depress. & Avg \\
\midrule
Raw Acc.\ (\%)     & 59.6 & 60.2 & 59.1 & 65.0 & 38.2 & 43.6 & 54.3 \\
Aligned Acc.\ (\%) & 66.3 & 56.5 & 57.1 & 72.3 & 61.5 & 54.3 & 61.3 \\
$\Delta$ (pp)      & \textbf{+6.7} & $-$3.6 & $-$2.0 & \textbf{+7.3} & \textbf{+23.3} & \textbf{+10.7} & \textbf{+7.0} \\
\midrule
Raw $\kappa$       & 0.195 & 0.214 & 0.182 & 0.298 & $-$0.232 & $-$0.128 & 0.088 \\
Aligned $\kappa$   & 0.296 & 0.129 & 0.154 & 0.440 & 0.239 & 0.088 & 0.224 \\
$n$ (pool pairs)   & 83 & 69 & 77 & 83 & 65 & 70 & --- \\
\bottomrule
\end{tabular}
\end{table}

Notably, only Mode~4 benefits substantially from LWRR on \emph{wealthy} (+7.3\,pp).
Modes~1 and 3 show \emph{negative} calibration gains ($-$10.1 and $-$11.6\,pp respectively), indicating that single-shot and single-image scores are too noisy for LWRR to learn meaningful local mappings.
Mode~2 breaks even ($\pm$0.0\,pp), suggesting that pairwise context alone, without multi-agent deliberation, produces scores that are marginally calibrated but not improved by LWRR.

Across all six categories (\cref{tab:vrm_gain_allcat}), Mode~4 LWRR yields a +7.0\,pp average improvement.
The most dramatic gains occur on categories where raw VLM performance is near chance: \emph{boring} (+23.3\,pp, from 38.2\% to 61.5\%) and \emph{depressing} (+10.7\,pp, from 43.6\% to 54.3\%).
Two categories show small negative gains (\emph{beautiful}: $-$3.6\,pp; \emph{lively}: $-$2.0\,pp), suggesting that for aesthetics-related perception, the VLM's raw dimension scores may already be near-optimally calibrated and LWRR introduces slight overfitting with the given reference set size.


% ==================================================================
% G. Complete Algorithm Pseudocode
% ==================================================================
\section{Complete Algorithm Pseudocode}
\label{sec:algorithm}

\begin{algorithm}[ht]
\caption{UrbanAlign: Post-hoc VLM Calibration for Urban Perception}
\label{alg:urbanalign}
\small
\begin{algorithmic}[1]
\REQUIRE $\mathcal{D}_{L}$: labelled pairs; $\phi_{\text{CLIP}}$: CLIP encoder; VLM; $K,\tau,\lambda,\alpha,\varepsilon,\theta$
\ENSURE $\mathcal{D}_{\mathrm{aligned}}$: calibrated dataset
\STATE \textbf{// Stage 1: Semantic Dimension Extraction}
\STATE Compute TrueSkill ratings $\{(\mu_i,\sigma_i)\}$ from $\mathcal{D}_{L}$
\STATE Sample $\mathcal{S}_{\mathrm{high}},\mathcal{S}_{\mathrm{low}}$ via consensus criteria
\STATE $\mathcal{D}_c \leftarrow \text{VLM}(\mathcal{S}_{\mathrm{high}}\cup\mathcal{S}_{\mathrm{low}})$ \COMMENT{$n$ dimensions}
\STATE \textbf{// Stage 2: Multi-Agent Feature Distillation}
\STATE Sample $\mathcal{D}_{\mathrm{sample}}\subset\mathcal{D}_{L}$; split into $\mathcal{D}_{\mathrm{ref}}, \mathcal{D}_{\mathrm{pool}}$
\FOR{each image pair $(x_A,x_B)\in\mathcal{D}_{\mathrm{sample}}$}
  \STATE $\text{obs}\leftarrow\text{VLM}_{\mathrm{Obs}}(x_A,x_B,\mathcal{D}_c)$
  \STATE $\text{debate}\leftarrow\text{VLM}_{\mathrm{Deb}}(x_A,x_B,\text{obs},\mathcal{D}_c)$
  \STATE $S(x_A),S(x_B)\leftarrow\text{VLM}_{\mathrm{Judge}}(x_A,x_B,\text{obs},\text{debate},\mathcal{D}_c)$
\ENDFOR
\STATE \textbf{// Stage 3: LWRR Calibration}
\STATE Build reference manifold $\mathcal{R}$ from $\mathcal{D}_{\mathrm{ref}}$ with mirror augment
\FOR{each $(x_A,x_B)\in\mathcal{D}_{\mathrm{pool}}$}
  \STATE Compute $\Delta_q^{\mathrm{hybrid}}$; find $K$-NN $\mathcal{N}_K$ in $\mathcal{R}$
  \STATE Solve LWRR (Eq.~(4) in the main paper); compute $\hat{\delta}, R^{2}$
  \STATE Re-infer $\hat{y}$ via Eq.~(6) in the main paper
\ENDFOR
\RETURN $\mathcal{D}_{\mathrm{aligned}}$
\end{algorithmic}
\end{algorithm}


% ---- Bibliography ----
\bibliographystyle{splncs04}
\bibliography{main}

\end{document}
